\documentclass{article}
\usepackage{amsmath,amssymb,amsthm}
\usepackage{algorithm}
\usepackage{algorithmic}
\usepackage{hyperref}
\usepackage{enumitem}
\newtheorem{theorem}{Theorem}
\newtheorem{lemma}{Lemma}
\newtheorem{assumption}{Assumption}
\newcommand{\E}{\mathbb{E}}
\newcommand{\R}{\mathbb{R}}

\begin{document}

\section*{Gradient Descent-Ascent (GDA)}

\section*{Mathematical Formulation}
Let $L:\R^{n}\times\R^{m}\rightarrow\R$ be a differentiable \emph{non‑convex–concave} objective.  
We wish to solve
\[
\min_{x\in\R^{n}}\;\max_{y\in\R^{m}} L(x,y).
\]
Define the stacked variable $z:=(x,y)\in\R^{n+m}$ and the \emph{gradient operator}
\[
\mathcal{G}(z)\;:=\;\begin{bmatrix}
-\nabla_{x}L(x,y)\\[2mm]
\;\;\nabla_{y}L(x,y)
\end{bmatrix}.
\]
GDA performs a simultaneous gradient \emph{ascent} step on $x$ and \emph{descent} step on $y$ with a common stepsize $\eta>0$:
\[
\boxed{
\begin{aligned}
x_{t+1} &= x_{t} + \eta\,\nabla_{x}L(x_{t},y_{t}),\\
y_{t+1} &= y_{t} - \eta\,\nabla_{y}L(x_{t},y_{t}).
\end{aligned}}
\tag{1}
\]

\section*{Pseudocode}
\begin{algorithm}[H]
\caption{Gradient Descent‑Ascent (GDA)}\label{alg:GDA}
\begin{algorithmic}[1]
\REQUIRE Initial point $(x_{0},y_{0})$, stepsize $\eta>0$, maximum iterations $T$.
\FOR{$t=0,\dots,T-1$}
    \STATE Compute stochastic (or exact) gradients 
        $g_{x}^{t}= \nabla_{x}L(x_{t},y_{t})$, 
        $g_{y}^{t}= \nabla_{y}L(x_{t},y_{t})$.
    \STATE Update 
        $x_{t+1}=x_{t}+ \eta\, g_{x}^{t}$,
        $y_{t+1}=y_{t}- \eta\, g_{y}^{t}$.
\ENDFOR
\RETURN $(x_{T},y_{T})$.
\end{algorithmic}
\end{algorithm}

\section*{Dry Run Example}
Consider the scalar function $f(w)=(w-3)^{2}$ (here $x\equiv w$, $y$ is absent).  
We apply (1) with stepsize $\eta=0.1$ and $x_{0}=0$.

\begin{center}
\begin{tabular}{c|c|c|c}
$t$ & $g_{x}^{t}=\nabla f(x_{t})$ & $x_{t+1}=x_{t}+ \eta g_{x}^{t}$ & $f(x_{t+1})$\\\hline
0 & $2(x_{0}-3)= -6$ & $0+0.1(-6)= -0.6$ & $( -0.6-3)^{2}=12.96$\\
1 & $2(-0.6-3)= -7.2$ & $-0.6+0.1(-7.2)= -1.32$ & $(-1.32-3)^{2}=18.6624$\\
2 & $2(-1.32-3)= -8.64$ & $-1.32+0.1(-8.64)= -2.184$ & $(-2.184-3)^{2}=26.8169$\\
\end{tabular}
\end{center}
The table shows the gradient evaluation, the ascent (here descent because $L=-f$) update, and the resulting objective value after three iterations.

\section*{Assumptions}
\begin{assumption}[Smoothness]\label{ass:smooth}
$L$ has $L$‑Lipschitz continuous gradients, i.e.
\[
\|\nabla L(z)-\nabla L(z')\|\le L\|z-z'\|,\qquad\forall z,z'\in\R^{n+m}.
\]
\end{assumption}
\begin{assumption}[Bounded Gradient]\label{ass:bound}
There exists $G>0$ such that $\|\nabla L(z)\|\le G$ for all $z$.
\end{assumption}
\begin{assumption}[Unbiased Stochastic Gradients]\label{ass:unbiased}
If stochastic gradients $g^{t}$ are used, then
\[
\E\big[\,g^{t}\mid z_{t}\,\big]=\nabla L(z_{t}),\qquad 
\E\big[\|g^{t}-\nabla L(z_{t})\|^{2}\mid z_{t}\big]\le\sigma^{2}.
\]
\end{assumption}
All constants $L,G,\sigma$ are finite and known.

\section*{Main Convergence Theorem}
\begin{theorem}\label{thm:main}
Assume \ref{ass:smooth}–\ref{ass:unbiased} and run GDA with a constant stepsize $\eta\le \frac{1}{2L}$.  
Define the \emph{stationarity measure}
\[
\Phi_{t}:=\|\nabla_{x}L(x_{t},y_{t})\|^{2}+\|\nabla_{y}L(x_{t},y_{t})\|^{2}.
\]
Then for any $T\ge 1$,
\[
\frac{1}{T}\sum_{t=0}^{T-1}\E[\Phi_{t}]
\;\le\;
\frac{2\big(L(x_{0},y_{0})-L^{*}\big)}{\eta T}+2\eta L\sigma^{2},
\tag{2}
\]
where $L^{*}:=\min_{x}\max_{y}L(x,y)$.  
Choosing $\eta =\Theta\!\big(\tfrac{1}{\sqrt{T}}\big)$ yields
\[
\min_{0\le t<T}\E[\Phi_{t}] = O\!\big(T^{-1/2}\big).
\]
\end{theorem}

\section*{Theoretical Properties}
\begin{itemize}[leftmargin=*]
\item \textbf{Stability.} The condition $\eta\le \frac{1}{2L}$ guarantees that the update (1) is a non‑expansive map for the gradient operator $\mathcal{G}$, preventing divergence even when $L$ is non‑convex in $x$.
\item \textbf{Hyperparameter Sensitivity.} The bound (2) shows a trade‑off: a larger $\eta$ speeds up the $1/(\eta T)$ term but inflates the variance term $2\eta L\sigma^{2}$. The optimal $\eta\asymp 1/\sqrt{T}$ balances both.
\item \textbf{Comparison to Extragradient.} Extragradient methods attain $O(1/T)$ under monotonicity. In the non‑convex–concave regime GDA only guarantees $O(1/\sqrt{T})$, matching known lower bounds for first‑order methods.
\end{itemize}

\section*{Discussion}
The theorem establishes that GDA converges to a \emph{first‑order stationary point} of the primal–dual gradient operator at the optimal rate $O(T^{-1/2})$. The result holds without any convexity/concavity assumptions beyond smoothness, making it applicable to a broad class of modern adversarial learning problems (e.g., GANs).

\section*{Preliminaries (Definitions and Lemmas)}
\begin{lemma}[Descent Lemma for Smooth Functions]\label{lem:descent}
If $\nabla L$ is $L$‑Lipschitz, then for any $z,z'\in\R^{n+m}$,
\[
L(z')\le L(z)+\langle\nabla L(z),z'-z\rangle+\frac{L}{2}\|z'-z\|^{2}.
\]
\end{lemma}
\begin{proof}
Standard proof by integrating the Lipschitz condition; omitted for brevity. \qedhere
\end{proof}

\section*{Main Proof (Step‑by‑Step Derivation)}
We prove Theorem~\ref{thm:main}. Throughout we denote $z_{t}=(x_{t},y_{t})$ and $g^{t}= (g_{x}^{t},-g_{y}^{t})$, where $g_{x}^{t}$ and $g_{y}^{t}$ are the (possibly stochastic) gradients used in \eqref{1}.  

\begin{enumerate}
\item[Step 1.] Apply Lemma~\ref{lem:descent} to $L$ at $z_{t}$ and $z_{t+1}=z_{t}+\eta g^{t}$:
\[
L(z_{t+1})\le L(z_{t})+\langle\nabla L(z_{t}),\eta g^{t}\rangle+\frac{L\eta^{2}}{2}\|g^{t}\|^{2}.
\tag{3}
\]

\item[Step 2.] Since $g^{t}$ is an unbiased estimator of $\nabla L(z_{t})$ (Assumption~\ref{ass:unbiased}),
\[
\E\big[\langle\nabla L(z_{t}),g^{t}\rangle\mid z_{t}\big]
= \|\nabla L(z_{t})\|^{2}.
\tag{4}
\]

\item[Step 3.] Take total expectation of \eqref{3} and use \eqref{4}:
\[
\E\big[L(z_{t+1})\big]\le \E\big[L(z_{t})\big]+\eta\E\big[\|\nabla L(z_{t})\|^{2}\big]
+\frac{L\eta^{2}}{2}\E\big[\|g^{t}\|^{2}\big].
\tag{5}
\]

\item[Step 4.] Decompose $\|g^{t}\|^{2}$ into variance and bias:
\[
\E\big[\|g^{t}\|^{2}\mid z_{t}\big]
= \|\nabla L(z_{t})\|^{2}+\E\big[\|g^{t}-\nabla L(z_{t})\|^{2}\mid z_{t}\big]
\le \|\nabla L(z_{t})\|^{2}+\sigma^{2},
\tag{6}
\]
where the inequality uses Assumption~\ref{ass:unbiased}.

\item[Step 5.] Substitute \eqref{6} into \eqref{5} and simplify:
\[
\E\big[L(z_{t+1})\big]\le \E\big[L(z_{t})\big]
+\eta\E\big[\|\nabla L(z_{t})\|^{2}\big]
+\frac{L\eta^{2}}{2}\Big(\E\big[\|\nabla L(z_{t})\|^{2}\big]+\sigma^{2}\Big).
\tag{7}
\]

\item[Step 6.] Regroup terms containing $\E[\|\nabla L(z_{t})\|^{2}]$:
\[
\E\big[L(z_{t+1})\big]\le \E\big[L(z_{t})\big]
+\Big(\eta+\frac{L\eta^{2}}{2}\Big)\E\big[\|\nabla L(z_{t})\|^{2}\big]
+\frac{L\eta^{2}}{2}\sigma^{2}.
\tag{8}
\]

\item[Step 7.] Choose $\eta\le\frac{1}{2L}$, which implies $L\eta\le\frac12$ and therefore
\[
\eta+\frac{L\eta^{2}}{2}\le \eta+\frac{\eta}{4}= \frac{5}{4}\eta
\le 2\eta.
\tag{9}
\]
Using (9) in (8) yields
\[
\E\big[L(z_{t+1})\big]\le \E\big[L(z_{t})\big]
+2\eta\,\E\big[\|\nabla L(z_{t})\|^{2}\big]
+\frac{L\eta^{2}}{2}\sigma^{2}.
\tag{10}
\]

\item[Step 8.] Rearrange (10) to isolate the gradient term:
\[
2\eta\,\E\big[\|\nabla L(z_{t})\|^{2}\big]
\ge \E\big[L(z_{t})-L(z_{t+1})\big]-\frac{L\eta^{2}}{2}\sigma^{2}.
\tag{11}
\]

\item[Step 9.] Sum \eqref{11} over $t=0,\dots,T-1$ and telescope the objective:
\[
2\eta\sum_{t=0}^{T-1}\E\big[\|\nabla L(z_{t})\|^{2}\big]
\ge \E\big[L(z_{0})-L(z_{T})\big]-\frac{L\eta^{2}\sigma^{2}T}{2}.
\tag{12}
\]

\item[Step 10.] Since $L(z_{T})\ge L^{*}$, we obtain
\[
2\eta\sum_{t=0}^{T-1}\E\big[\|\nabla L(z_{t})\|^{2}\big]
\ge L(z_{0})-L^{*}-\frac{L\eta^{2}\sigma^{2}T}{2}.
\tag{13}
\]

\item[Step 11.] Divide both sides of \eqref{13} by $2\eta T$:
\[
\frac{1}{T}\sum_{t=0}^{T-1}\E\big[\|\nabla L(z_{t})\|^{2}\big]
\le \frac{L(z_{0})-L^{*}}{\eta T}+\frac{L\eta\sigma^{2}}{2}.
\tag{14}
\]

\item[Step 12.] Recall that $\|\nabla L(z_{t})\|^{2}
= \|\nabla_{x}L(x_{t},y_{t})\|^{2}+\|\nabla_{y}L(x_{t},y_{t})\|^{2}
= \Phi_{t}$.  Hence (14) is exactly the bound (2) claimed in Theorem~\ref{thm:main}.
\qedhere
\end{enumerate}

\section*{Rate Derivation (With Constants)}
From \eqref{14},
\[
\frac{1}{T}\sum_{t=0}^{T-1}\E[\Phi_{t}]
\le \frac{\Delta_{0}}{\eta T}+ \frac{L\eta\sigma^{2}}{2},
\qquad \Delta_{0}:=L(z_{0})-L^{*}.
\]
Set $\eta = \sqrt{\frac{2\Delta_{0}}{L\sigma^{2}}}\,T^{-1/2}$, which respects $\eta\le \frac{1}{2L}$ for all sufficiently large $T$. Substituting,
\[
\frac{1}{T}\sum_{t=0}^{T-1}\E[\Phi_{t}]
\le 2\sqrt{\frac{L\Delta_{0}\sigma^{2}}{2}}\,T^{-1/2}
= O\!\big(T^{-1/2}\big).
\]
Consequently there exists at least one iterate $t$ with $\E[\Phi_{t}] = O(T^{-1/2})$.

\section*{Special Cases}
\begin{enumerate}[label=\arabic*)]
\item \textbf{Deterministic Gradients ($\sigma=0$).}  
Bound (2) reduces to $\frac{1}{T}\sum_{t=0}^{T-1}\E[\Phi_{t}]\le \frac{2\Delta_{0}}{\eta T}$.
Choosing a constant $\eta\le\frac{1}{2L}$ yields a $O(1/T)$ decay of the average stationarity measure.
\item \textbf{Pure Minimization ($y$ absent).}  
The algorithm collapses to standard gradient descent. The proof above simplifies to the classical $O(1/\sqrt{T})$ rate for non‑convex smooth objectives.
\item \textbf{Convex–Concave Setting.}  
If $L$ is convex in $x$ and concave in $y$, the same bound holds, but stronger results (e.g., $O(1/T)$ in primal–dual gap) can be obtained via extragradient or mirror‑prox methods.
\end{enumerate}

\end{document}